\documentclass[12pt,oneside,a4paper,article]{abntex2}
\usepackage[utf8]{inputenc}
\usepackage{array}
\usepackage[T1]{fontenc}
\usepackage[brazil]{babel}
\usepackage{graphicx}
\usepackage{tabularx}
\usepackage{amsmath}
\usepackage{authblk}
\usepackage{mathptmx}
\usepackage{xcolor}
\usepackage{float}
\usepackage{tikz}
\usetikzlibrary{positioning,arrows.meta,shapes.geometric,fit,calc}

\usepackage[left=3cm,right=2cm,top=3cm,bottom=2cm]{geometry}
\usepackage{setspace}
\linespread{1.25}

\title{\textbf{Testes e Qualidade de Software}\\[0.35em]
{\large Simulador de Batalha -- PokeSal}\\[0.25em]}

\author{Guilherme Monteiro Da Silva
\thanks{guilhermemonteiro.silva@ucsal.edu.br}}

\author{Miguel Novaes Palma Santos
\thanks{miguelnoaves.santos@ucsal.edu.br}}

\author{Dantton Lavigne
\thanks{dantton.pinheiro@ucsal.edu.br}}

\author{Guilherme Gomes de Magalhães
\thanks{guilherme.magalhaes@ucsal.edu.br}}

\author[1*]{Pedro Arthur de Melo Nascimento
\thanks{pedroam.nascimento@pro.ucsal.br}}

\affil{
Engenharia de Software \par
Escola de Tecnologias \par
Universidade Católica do Salvador (UCSAL) \par
Av. Prof. Pinto de Aguiar, 2589 Pituaçu, CEP: 41740-090 \par
Salvador/BA, Brasil
}

\date{Setembro de 2026}

\ifthenelse{\equal{\ABNTEXisarticle}{true}}{%
\renewcommand{\maketitlehookb}{}
}{}

\usepackage{hyperref}
\hypersetup{
colorlinks=false,
pdfborder={0 0 0},
}

\renewcommand*{\Authsep}{, }
\renewcommand*{\Authand}{, }
\renewcommand*{\Authands}{, }
\renewcommand*{\Affilfont}{\normalsize\normalfont}
\renewcommand*{\Authfont}{\bfseries}
\setlength{\affilsep}{2em}
\newsavebox\affbox

\begin{document}

\begin{center}
\IfFileExists{imagens-template/ucsal_logo.png}{%
\includegraphics[width=0.3\textwidth]{imagens-template/ucsal_logo.png}%
}{}
\end{center}
{\let\newpage\relax\maketitle}

\begin{resumoumacoluna}
Relatório da Fase 01 do PokeSal, da Equipe Chaos, na matéria de Testes e
Qualidade de Software. O jogo é uma batalha 1 contra 1 no estacionamento da
UCSal: cada um escolhe um inicial, o terreno mexe no dano e a luta vai até
alguém cair. Aqui entram a leitura do enunciado, os três requisitos autorais,
os diagramas UML, o que fizemos no código para seguir o Google Java Style,
os bugs que apareceram no caminho e o trabalho de cada um.
\noindent
\newline
\textbf{Palavras-chaves}: PokeSal. Requisitos. UML. Código limpo. Checkstyle.
\end{resumoumacoluna}

\textual

\tableofcontents
\newpage

\section{Introdução}

O trabalho é um torneio de batalhas no estacionamento da UCSal Pituaçu. Cada
treinador entra com um inicial (BulbaSal, CharSal, SquirtSal, ChikoSal,
CyndaSal ou TotoSal). A luta é por turno, com tipo, terreno, status e no
máximo dois itens.

Na Fase 01, entregamos a análise dos requisitos, três regras autorais,
UML (casos de uso e classes), código no Git e um relatório dizendo o que cada um fez.

\section{Análise estática de requisitos}
\label{sec:inspecao}

Lemos o enunciado antes de sair programando, para marcar o que estava
claro, o que estava solto e o que a gente ia ter que decidir.

\subsection{O que já vinha pronto}

\begin{itemize}
  \item um inicial só, entre os seis nomes;
  \item HP, ATK, DEF, SPD e tipo;
  \item Fogo, Água e Planta ($2{,}0$x e $0{,}5$x);
  \item Asfalto Quente: Fogo $+15\%$ de dano;
  \item Poça de Chuva: Golpes de Água aplicam 10\% adicionais de precisão ou dano;
  \item Canteiro Central: Planta recupera $5\%$ do HP máximo no fim do turno;
  \item ordem do turno pela SPD;
  \item status no fim do turno (Queimado, Envenenado, Paralisado);
  \item no máximo dois itens, e item gasta o turno.
\end{itemize}

\subsection{O que o texto deixou em aberto}

A poça de chuva dá ``$10\%$ a mais de precisão ou dano'' para Água. Precisão e
dano não são a mesma coisa. Se a gente colocasse chance de errar, cada teste
dava um resultado diferente. Ficamos com \textbf{dano}.

O enunciado cita os status, mas não diz quem aplica. Não tem item de status
nem sorteio. A gente ligou no golpe elemental: Fogo queima, Água paralisa,
Planta envenena.

Queimado ``reduz HP e ATK'' sem número. Ficou metade do ATK e $10\%$ do HP
máximo por turno. Veneno ``progressivo'' também sem detalhe: percentual do
HP máximo subindo a cada turno ($5\%$, $10\%$, $15\%$\ldots). Paralisia
corta a SPD pela metade.

Fora isso, o guia do projeto não fala de stats dos iniciais, fórmula de dano, empate de
SPD, limite de turno, se pode empatar, se status acumula, quanto a Potion
cura, o que acontece se o HP passar de zero ou do máximo, nem o que fazer
com nome fora da lista.

A fórmula a gente fechou só para o golpe ter um número. ``Um status por vez'' e ``lutar até o fim'' a gente escolheu de
proposito, porque o enunciado não cobre.

Quase-contradição: o texto manda ordenar \textbf{estritamente} por SPD e não
fala de empate. SPD igual precisa de alguma regra. Caso haja empate de SPD, começa quem tem mais ATK; se o ATK empatar também, o treinador 1 começa.

\section{Requisitos autorais}
\label{sec:autorais}

Três regras nossas, já no sistema e com testes feitos. Desempate de SPD
por ATK \textbf{não entra} aqui.

\subsection{RA01 -- Combo}

O mesmo golpe duas vezes seguidas: o segundo ganha $+50\%$. Tem golpe
elemental (tipo + terreno) e Investida (neutra). Trocar de golpe ou usar item
zera o combo. A Investida existe também para dar para testar o reset: se só
tivesse um golpe, não tinha como mostrar ``o outro reseta''.

\subsection{RA02 -- Lutar até o fim}

A luta só acaba quando um cai. Sem teto de turno e sem empate. O caso chato
é os dois zerarem no mesmo tick de status. Aí a gente olha o HP de cada um
logo antes desse tick: ganha quem tinha mais vida. Se era igual, o treinador 1 vence.

\subsection{RA03 -- Um status por vez}

Não junta Queimado + Veneno + Paralisado. O novo troca o antigo. Antidote
limpa. Se o veneno entra no lugar da queimadura, o contador do veneno volta
para zero. O enunciado lista os três e não diz se somam; somar deixava o fim
de turno uma bagunça.

\section{Modelagem UML}
\label{sec:uml}

O sistema é todo feito no console (CLI), 1 contra 1, um inicial para cada um, um terreno.

\subsection{Diagrama de casos de uso}

O ator é o treinador (os dois fazem as mesmas coisas): escolher inicial,
escolher terreno, golpe elemental, Investida, usar item e encerrar a batalha
quando alguém cai.

\begin{figure}[H]
\centering
\begin{tikzpicture}[
  font=\small,
  ator/.style={circle, draw, minimum size=0.55cm},
  uc/.style={ellipse, draw, minimum width=3.3cm, minimum height=0.95cm,
    align=center, inner sep=2pt},
  sistema/.style={draw, rounded corners, inner sep=10pt}
]
  \node[ator] (cabeca) {};
  \draw ($(cabeca.south)+(0,-0.05)$) -- ++(0,-0.85);
  \draw ($(cabeca.south)+(0,-0.35)$) -- ++(-0.35,-0.25);
  \draw ($(cabeca.south)+(0,-0.35)$) -- ++(0.35,-0.25);
  \draw ($(cabeca.south)+(0,-0.90)$) -- ++(-0.3,-0.45);
  \draw ($(cabeca.south)+(0,-0.90)$) -- ++(0.3,-0.45);
  \node[below] at ($(cabeca.south)+(0,-1.45)$) {Treinador};

  \node[uc, right=2.4cm of cabeca, yshift=2.3cm] (c1) {Escolher inicial};
  \node[uc, below=0.28cm of c1] (c2) {Escolher terreno};
  \node[uc, below=0.28cm of c2] (c3) {Golpe elemental};
  \node[uc, below=0.28cm of c3] (c4) {Usar Investida};
  \node[uc, below=0.28cm of c4] (c5) {Usar item};
  \node[uc, below=0.28cm of c5] (c6) {Encerrar batalha};

  \node[sistema, fit=(c1)(c2)(c3)(c4)(c5)(c6),
    label=above:Sistema PokeSal] {};

  \foreach \n in {c1,c2,c3,c4,c5,c6}
    \draw[-{Stealth}] (cabeca.east) -- (\n.west);
\end{tikzpicture}
\caption{Casos de uso do PokeSal.}
\label{fig:casos-uso}
\end{figure}

Inicial fora da lista não entra. Terceiro item também não. Encerrar batalha
não é botão do menu: acontece quando o HP chega a zero, e sempre tem
vencedor (RA02).

\subsection{Diagrama de classes}

Domínio em \texttt{pokesal.modelo}, luta em \texttt{pokesal.batalha},
exceção em \texttt{pokesal.excecao}. O \texttt{Main} só lê o que o usuário
digitou e manda a batalha resolver o turno.

\begin{figure}[H]
\centering
\resizebox{\textwidth}{!}{%
\begin{tikzpicture}[
  font=\scriptsize,
  cls/.style={draw, rectangle, align=left, inner sep=4pt},
  enu/.style={draw, rectangle, align=left, inner sep=4pt, fill=gray!8}
]
  \node[cls] (treinador) {%
    \textbf{Treinador}\\[-1pt]
    \hline
    $-$ nome: String\\
    $-$ pokesal: Pokesal\\
    $-$ itensUsados: int\\[-1pt]
    \hline
    $+$ usarItem(item)\\
    $+$ getItensRestantes()
  };

  \node[cls, right=0.55cm of treinador] (pokesal) {%
    \textbf{Pokesal}\\[-1pt]
    \hline
    $-$ nome, tipo, hp, atk, def, spd\\
    $-$ status: StatusEfeito\\
    $-$ ultimoGolpe: TipoGolpe\\[-1pt]
    \hline
    $+$ getAtkEfetivo() / getSpdEfetivo()\\
    $+$ aplicarStatus(status)\\
    $+$ aplicarEfeitosStatusFimTurno()\\
    $+$ isCombo(golpe) / registrarGolpe()
  };

  \node[cls, below=0.7cm of treinador] (batalha) {%
    \textbf{Batalha}\\[-1pt]
    \hline
    $-$ treinador1, treinador2\\
    $-$ terreno: Terreno\\[-1pt]
    \hline
    $+$ obterOrdemDeAtaque()\\
    $+$ executarAcao(...)\\
    $+$ resolverTurno(...)\\
    $+$ obterVencedor()
  };

  \node[cls, right=0.55cm of batalha] (calc) {%
    \textbf{CalculadoraDano}\\[-1pt]
    \hline
    \hline
    $+$ calcular(...)\\
    $+$ obterMultiplicadorTerreno(...)
  };

  \node[cls, below=0.7cm of batalha] (cat) {%
    \textbf{CatalogoInicial}\\[-1pt]
    \hline
    \hline
    $+$ criar(nome)\\
    $+$ nomesPermitidos()
  };

  \node[enu, right=0.55cm of cat] (enums) {%
    \textbf{Enumerações}\\[-1pt]
    \hline
    TipoElemental \quad Terreno\\
    StatusEfeito \quad TipoItem\\
    AcaoTurno \quad TipoGolpe
  };

  \node[cls, below=0.7cm of cat] (ex) {%
    \textbf{Exceções}\\[-1pt]
    \hline
    InicialInvalidoException\\
    LimiteItensExcedidoException
  };

  \draw[-{Stealth}] (treinador) -- node[above, font=\tiny] {1} (pokesal);
  \draw[-{Stealth}] (batalha.north) -- (treinador.south);
  \draw[-{Stealth}] (batalha) -- (calc);
  \draw[-{Stealth}] (batalha.east) -- ++(0.35,0) |- (pokesal.south);
\end{tikzpicture}%
}
\caption{Classes da Fase 01.}
\label{fig:classes}
\end{figure}

\texttt{Treinador} tem um \texttt{Pokesal} e conta item.
\texttt{Batalha} guarda os dois, o terreno, a ordem por SPD e o fim do turno.
\texttt{Pokesal} segura HP, status e combo. Nome inválido vira
\texttt{InicialInvalidoException}. Terceiro item vira
\texttt{LimiteItensExcedidoException}.

\section{Implementação e boas práticas (Checkstyle)}
\label{sec:checkstyle}

O código está no repositório Git, nos pacotes \texttt{pokesal}, \texttt{modelo},
\texttt{batalha} e \texttt{excecao}. A matéria pede Google Java Style. O que
a gente aplicou:

\begin{itemize}
  \item classe em PascalCase, método/variável em camelCase, constante em
        UPPER\_SNAKE;
  \item Javadoc em classe e método público;
  \item número fixo de dano, vantagem, item e status em
        \texttt{ConstantesBatalha};
  \item linha por volta de 100 caracteres e indentação de 2 espaços;
  \item import um a um, sem importar o pacote inteiro ou incluir WildCards.
\end{itemize}

O console pergunta treinador, inicial e terreno, e roda turno até ter
vencedor.

\section{Problemas encontrados e correções}
\label{sec:bugs}

Isso aqui apareceu quando a gente foi montar e testar no console:

\textbf{Null na hora de criar a luta.} Dava para subir batalha sem treinador,
sem Pokésal ou sem terreno. Quebrava no meio do turno, com erro genérico.
Agora isso é recusado na entrada.

\textbf{HP negativo e cura sem teto.} Golpe forte ia para baixo de zero;
Potion com a vida cheia passava do máximo. HP para em zero; cura para no
máximo da espécie.

\textbf{Inicial inventado.} Qualquer nome passava. Agora só os seis; o resto
cai numa exceção só para isso.

\textbf{Terceiro item.} O limite era dois, mas a terceira tentativa não
falhava de um jeito que o teste pegasse. A gente passou a lançar a exceção de limite
e, no console, a pessoa escolhe outra ação em vez de o programa morrer.

\textbf{SPD igual.} Os dois com a mesma velocidade: a ordem mudava sem
regra. Como o enunciado não cobre, desempate por ATK (não é autoral).

\textbf{Os dois caíam no status.} Dois queimados, pouca vida, o tick zerava
os dois e não tinha vencedor. A gente guarda o HP de antes do tick e usa
isso para decidir.

\textbf{Combo depois do item.} O $+50\%$ seguia mesmo depois de usar item.
Item gasta o turno, então o combo zera.

\textbf{Status empilhado.} Dava para estar queimado e envenenado juntos, e o
fim do turno cobrava os dois. Virou o RA03: o novo troca o antigo.

\textbf{Número solto e import de pacote.} Percentual escrito na conta e
import com asterisco. Foi para constante nomeada e import explícito, que é
o que o Checkstyle pede.

Se a gente entregasse assim, isso ia aparecer na arguição. Por isso
corrigimos ainda nesta fase.

\section{Relatório de contribuição individual}
\label{sec:contribuicao}

A divisão saiu na reunião de 18/09/2026, 18h00--20h30, no Google Meet. A
segunda foi em 20/09/2026, 21h00--00h00: fechar código, revisar autoral,
corrigir o que estava quebrando e organizar o Git. Os quatro estavam nas
duas.

\subsection{Guilherme Monteiro Da Silva}

A inspeção de requisitos foi de Guilherme Monteiro. Ele leu o enunciado
antes da gente programar e marcou detalhes importantes, como a observação da poça de chuva deixar em aberto ``precisão ou dano'', falta de
fórmula, empate de SPD, status sem número e o fato de o texto não dizer
como o status entra na luta. Isso virou a Seção~\ref{sec:inspecao}.

\subsection{Guilherme Gomes de Magalhães}

Os três autorais são de Guilherme Gomes: Combo, lutar até o fim e um status
por vez. Também ajudou a deixar claro que desempate de SPD por ATK não
é um requisito autoral, e sim correção do enunciado. O texto da Seção~\ref{sec:autorais} é
desse trabalho.

\subsection{Miguel Novaes Palma Santos}

Os diagramas UML são de Miguel: casos de uso e classes, batendo com o
que o PokeSal faz de verdade (inicial, terreno, golpe, item, encerrar;
\texttt{Treinador}, \texttt{Pokesal}, \texttt{Batalha},
\texttt{CalculadoraDano}, catálogo, enum e exceção).

\subsection{Dantton Lavigne}

Quem cuidou da maior parte do código foi Dantton: pacotes, interface no console, constante
nomeada, Javadoc, indentação, imports e o ambiente para o
resto da equipe rodar testes e apontar erros e acertos. Realizou também a primeira leva de testes com os nomes da
Fase 02 e a conferência do Checkstyle.

\subsection{O que a gente fez junto}

Os quatro rodaram o projeto, bateram os autorais com o que estava
funcionando, corrigiram o que está na Seção~\ref{sec:bugs} e conferiram o
que ia para o repositório.

\section{Uso de inteligência artificial}
\label{sec:ia}

A gente usou IA em duas frentes nesta fase: passar o olho na ortografia
deste relatório e entender erro que o compilador apontava. Não foi a
ferramenta escrevendo o trabalho no nosso lugar. A gente redigia, compilava,
e quando o texto ou o build quebrava pedia ajuda para achar o que estava
errado.

No relatório, o que mais aparecia era acento, concordância e frase pela
metade depois de um recorte. Um exemplo bobo: ``proposito'' sem acento.
Outro: a gente tinha combinado escrever ``de fulano'' (e não ``do fulano''),
e a revisão só checava se isso tinha escapado. A sugestão vinha, a gente
lia e aceitava ou não.

No compilador, dois tipos de erro repetiram. No Java, o Ctrl+F5 caía em
\texttt{ClassNotFoundException} no \texttt{Main}: o programa procurava a
classe em \texttt{target/classes} e a pasta estava vazia, porque o código
ainda não tinha sido compilado nessa pasta. No \LaTeX{}, o que mais
atrapalhou foi arquivo de logo que não existia
(\texttt{imagens-template/ucsal\_logo.png}), porcentagem solta no texto
(o compilador trata \texttt{\%} como comentário e come o resto da linha)
e nome de classe com underline fora de \texttt{texttt}, que o pdflatex
reclama. A IA explicava a mensagem; quem mexia no arquivo era a equipe.

Esses usos estão documentados no arquivo \texttt{AI\_DECLARATION.md}, na
raiz do repositório, junto com os prompts que a gente de fato usou.

\section{Considerações finais}

A Fase 01 termina com o jogo rodando, as regras do enunciado, três autorais, UML
batendo com o código e fonte no padrão de nome, Javadoc, constantes e imports.

O enunciado tem alguns buracos. A gente marcou, separou o que era correção e o que
era autoral, e consertou o que só aparece quando o programa roda (null, HP
inválido, item a mais, combo que não zerava, os dois caindo no mesmo tick).
Ata, Sonar, checklist de código e o documento de testes da Fase 02 ficam
para as próximas entregas.

\end{document}
